%% ============================================================
%% SECTION 7: HALLUCINATION IN RS-MLLMs
%% EarthVision-LM Survey -- Artificial Intelligence Review
%% Template: sn-jnl (Springer Nature) | sn-mathphys-num
%% W3 FIX: All three hallucination types literature-grounded
%%          Type III grounded in HM-Bench 2026 empirically
%% Rules: booktabs, TikZ, numbered [N] cites, NO fabricated data
%% ============================================================

\section{Hallucination in RS-MLLMs}\label{sec7}

Hallucination in RS-MLLMs is not merely an accuracy problem---it is a safety
problem. When an RS-MLLM fabricates a building where none exists, misidentifies
a military installation, or describes an empty harbour as containing vessels, the
consequences in operational contexts can be catastrophic. This section provides
a systematic hallucination taxonomy for RS-MLLMs that proposes a formal three-way
partition (Localization, Discrimination, and Spectral failure modes) grounded in
the RS domain challenges of Sect.~\ref{sec2}---extending prior work through an
empirically grounded third type. To the best of our knowledge, prior RS-MLLM
hallucination studies have not explicitly isolated spectral hallucination as a
distinct failure category; our taxonomy extends prior localisation/discrimination
frameworks by incorporating spectral failure evidence from HM-Bench~\cite{hmbench2026}.

%% ----------------------------------------------------------
\subsection{Why RS Hallucination Is More Dangerous}\label{sec7:danger}
%% ----------------------------------------------------------

Hallucination in general-domain MLLMs---describing objects absent from a
natural photograph---is primarily a reliability inconvenience. In remote sensing,
the same failure mode carries qualitatively higher stakes across three critical
application domains.

\textbf{Disaster response.}
Post-disaster RS-MLLM outputs guide search-and-rescue operations: teams are
directed to locations identified as containing collapsed structures, blocked
roads, or survivor clusters. A localization hallucination that places a
collapsed building 500\,m from its true position redirects rescue teams away
from survivors. In the 2023 Turkey earthquake response, RS-based damage mapping
errors were documented as a primary cause of delayed rescue in rural
areas~\cite{long2021creating}. The time cost of an RS-MLLM hallucination is
not a wrong caption---it is measured in lives.

\textbf{Infrastructure monitoring.}
Automated inspection of pipelines, power lines, and bridge structures via
RS-MLLMs depends on accurate damage classification. A discrimination
hallucination that misclassifies a structural crack as surface discolouration
triggers either dangerous neglect or unnecessary and costly mobilisation of
maintenance teams. Either failure mode has direct economic and safety
consequences disproportionate to those of a mislabelled natural image.

\textbf{Agricultural and environmental mapping.}
Crop-type misclassification propagates into yield estimates, insurance payouts,
and subsidy allocations affecting millions of smallholder farmers. A spectral
hallucination that interprets drought-stressed vegetation as bare soil---because
their RGB appearances are similar but their spectral signatures diverge in
the NIR---leads to systematically wrong irrigation and intervention decisions.

These high-stakes failure modes have no close parallel in the natural-image
MLLM setting, where the worst case of hallucination is typically a wrong
object label or a confabulated description. Existing work has begun to
acknowledge RS hallucination as a distinct research problem:
DDFAV~\cite{ddfav2025} identifies low-quality instruction data as a root
cause; RSHallu~\cite{rshallu2026} provides the first diagnostic benchmark;
Seeing Clearly~\cite{seeingclearly2026} proposes training-free mitigation.
Figure~\ref{fig:hall_example} illustrates a canonical RS hallucination failure.

\begin{figure}[h]
\centering
\begin{tikzpicture}[font=\small, >=Stealth, scale=0.78]
  %% Left panel: RS harbour scene
  \fill[gray!18] (0,0) rectangle (4.8,3.5);
  \fill[gray!25] (0.1,0.1) rectangle (4.7,3.4);
  %% Water (dark)
  \fill[blue!20!gray!22] (0.2,0.2) rectangle (4.6,2.2);
  %% Dock structure
  \fill[gray!55] (0.2,2.0) rectangle (2.2,2.35);
  \fill[gray!55] (0.2,2.2) rectangle (0.65,3.2);
  %% Second dock arm
  \fill[gray!50] (2.6,2.2) rectangle (4.5,2.5);
  %% Empty berths -- NO ships
  \fill[blue!18!gray!20] (0.7,2.15) rectangle (2.0,3.15);
  \fill[blue!18!gray!20] (2.1,2.15) rectangle (3.4,3.15);
  %% Berth labels
  \node[gray!55, font=\tiny] at (1.35,2.65) {Berth 1};
  \node[gray!55, font=\tiny] at (2.75,2.65) {Berth 2};
  %% Image label
  \node[font=\scriptsize\bfseries] at (2.45,0.55)
    {RS Image: Harbour};
  \node[font=\tiny, gray!65] at (2.45,0.25)
    {(Both berths empty)};

  %% Right panel: query + MLLM response
  \fill[gray!8] (5.2,0) rectangle (11.5,3.5);
  \draw[gray!40] (5.2,0) rectangle (11.5,3.5);
  %% Query box
  \fill[blue!10] (5.4,2.85) rectangle (11.3,3.35);
  \node[anchor=west, font=\scriptsize\bfseries] at (5.5,3.10)
    {Query: ``Are there ships in the harbour?''};
  %% MLLM hallucinated response
  \fill[red!10] (5.4,1.60) rectangle (11.3,2.80);
  \node[anchor=north west, font=\scriptsize, align=left,
        text width=5.5cm] at (5.5,2.78)
    {\textbf{MLLM Output:} ``Yes, I can clearly see
    two large cargo vessels docked in the berths.
    The ship on the left appears to be a bulk carrier
    approximately 150\,m in length.''};
  %% Ground truth
  \fill[green!10] (5.4,0.55) rectangle (11.3,1.55);
  \node[anchor=north west, font=\scriptsize, align=left,
        text width=5.5cm] at (5.5,1.53)
    {\textbf{Ground Truth:} Both berths are empty.
    No vessels present in the scene at acquisition time.};
  %% Labels
  \node[red!70, font=\tiny\bfseries] at (8.35,2.82)
    {HALLUCINATION};
  \node[green!50!black, font=\tiny\bfseries] at (8.35,1.57)
    {CORRECT};
  %% Type annotation
  \node[orange!80, font=\scriptsize\bfseries] at (8.35,0.28)
    {Failure Type: Type~I (Localization Hallucination)};
\end{tikzpicture}
\caption{Canonical RS hallucination failure (Type~I). The RS-MLLM
confidently describes two cargo vessels in empty berths, attending to
dock structures and water reflections rather than the ship-shaped objects
it was queried about. This failure pattern is consistent with the
``Cannot Find'' category documented in Aerial Mirage~\cite{aerialmirage2026}
and tested in the RSHallu benchmark~\cite{rshallu2026}.}\label{fig:hall_example}
\end{figure}

%% ----------------------------------------------------------
\subsection{Existing RS Hallucination Work}\label{sec7:prior}
%% ----------------------------------------------------------

We survey the four primary contributions to RS hallucination research before
presenting our taxonomy. We explicitly acknowledge where our taxonomy builds
on rather than replaces this prior work.

\textbf{DDFAV (MDPI Remote Sensing, February 2025)~\cite{ddfav2025}.}
DDFAV establishes that shallow, template-generated instruction data is a
primary structural cause of RS-MLLM hallucinations. The paper demonstrates
that models trained on low-quality instruction pairs over-fit to answer
templates rather than learning genuine spatial reasoning, producing confident
but physically ungrounded responses. DDFAV contributes a quality-controlled
instruction protocol and confirms that data quality outweighs data quantity
for reducing RS hallucination rates.

\textbf{Aerial Mirage (2026)~\cite{aerialmirage2026}.}
Aerial Mirage provides the first formal categorisation of RS hallucination
into two failure types based on the model's ability to localise the queried
target. \textit{Type~1 (``Cannot Find'')} occurs when model attention disperses
over irrelevant image regions and the target is never localised. \textit{Type~2
(``Cannot See Clearly'')} occurs when the correct region is attended to but
the fine-grained class is misidentified. This two-type framework directly
informs the first two types of our taxonomy; we adopt Aerial Mirage's
categorisation and extend it rather than proposing an independent alternative.

\textbf{RSHallu (2026)~\cite{rshallu2026}.}
RSHallu is a systematic benchmark designed as a diagnostic tool for
hallucinations in the remote sensing domain. It evaluates RS-MLLMs on
object existence, object attribute, and object relationship queries across
diverse RS scene types, documenting high confusion rates among visually
similar object classes (car\,/\,truck\,/\,bus in parking lots;
building\,/\,shadow\,/\,road in urban scenes). RSHallu provides the
primary empirical grounding for our Type~II (Discrimination) hallucination
category.

\textbf{Seeing Clearly (March 2026)~\cite{seeingclearly2026}.}
Seeing Clearly proposes a training-free inference-time mitigation for
RS hallucination by correcting diffuse attention maps before response
generation. The method addresses both Aerial Mirage types without
fine-tuning, confirming that hallucination in RS-MLLMs has an attentional
root cause amenable to inference-time correction.

\textbf{Relationship to our taxonomy.}
We extend Aerial Mirage's two-type framework with a third type
(Type~III: Spectral Hallucination) grounded in the empirical findings of
HM-Bench~\cite{hmbench2026} on spectral reasoning failure across 18 MLLMs.
Types~I and~II are adopted from Aerial Mirage with refined causal analysis;
Type~III is new and, to the best of our knowledge, prior RS-MLLM hallucination
studies have not explicitly isolated spectral hallucination as a distinct
failure category in the RS hallucination literature.

%% ----------------------------------------------------------
\subsection{A Three-Type RS-MLLM Hallucination Taxonomy}\label{sec7:taxonomy}
%% ----------------------------------------------------------

We propose a three-type taxonomy of RS-MLLM hallucination failures,
organised by the \emph{stage} at which the failure occurs in the model
processing pipeline (Fig.~\ref{fig:hall_tree}): failure to \emph{locate}
the query target (Type~I), failure to \emph{classify} a located target
(Type~II), and failure to \emph{interpret sensor modality} correctly
(Type~III).

\begin{figure}[h]
\centering
\begin{tikzpicture}[font=\small, >=Stealth,
    root/.style={draw, fill=gray!14, rounded corners=3pt,
                 minimum width=5.2cm, minimum height=0.68cm,
                 align=center},
    tI/.style={draw, fill=blue!14, rounded corners=3pt,
               minimum width=3.5cm, minimum height=0.68cm,
               align=center},
    tII/.style={draw, fill=orange!14, rounded corners=3pt,
                minimum width=3.5cm, minimum height=0.68cm,
                align=center},
    tIII/.style={draw, fill=red!14, rounded corners=3pt,
                 minimum width=3.5cm, minimum height=0.68cm,
                 align=center},
    sub/.style={draw=gray!45, fill=gray!6, rounded corners=2pt,
                minimum width=3.2cm, minimum height=0.52cm,
                align=center, font=\scriptsize},
    arr/.style={->, thick, gray!60, shorten >=1pt}]

  %% Root
  \node[root] (root) at (0,0)
    {RS-MLLM Hallucination};

  %% Three branches
  \node[tI]   (t1) at (-4.8,-1.6) {Type~I\\Localization};
  \node[tII]  (t2) at ( 0.0,-1.6) {Type~II\\Discrimination};
  \node[tIII] (t3) at ( 4.8,-1.6) {Type~III\\Spectral~[NEW]};

  \draw[arr] (root.south) -- ++(0,-0.3) -| (t1.north);
  \draw[arr] (root.south) -- (t2.north);
  \draw[arr] (root.south) -- ++(0,-0.3) -| (t3.north);

  %% Type I sub-nodes
  \node[sub] (t1a) at (-4.8,-2.7)  {``Cannot Find''};
  \node[sub] (t1b) at (-4.8,-3.4)  {Diffuse attention};
  \node[sub, fill=blue!8, draw=blue!40]
             (t1c) at (-4.8,-4.1)
    {\textit{Aerial Mirage} T1~\cite{aerialmirage2026}};
  \draw[arr] (t1)--(t1a); \draw[arr](t1a)--(t1b); \draw[arr](t1b)--(t1c);

  %% Type II sub-nodes
  \node[sub] (t2a) at ( 0.0,-2.7)  {``Cannot See Clearly''};
  \node[sub] (t2b) at ( 0.0,-3.4)  {Fine-grained confusion};
  \node[sub, fill=orange!8, draw=orange!40]
             (t2c) at ( 0.0,-4.1)
    {\textit{RSHallu}~\cite{rshallu2026} +
     \textit{Aerial Mirage} T2};
  \draw[arr] (t2)--(t2a); \draw[arr](t2a)--(t2b); \draw[arr](t2b)--(t2c);

  %% Type III sub-nodes
  \node[sub] (t3a) at ( 4.8,-2.7)  {RGB prior on non-RGB};
  \node[sub] (t3b) at ( 4.8,-3.4)  {Spectral physics failure};
  \node[sub, fill=red!8, draw=red!40]
             (t3c) at ( 4.8,-4.1)
    {\textit{HM-Bench} 2026~\cite{hmbench2026}};
  \draw[arr] (t3)--(t3a); \draw[arr](t3a)--(t3b); \draw[arr](t3b)--(t3c);

  %% Bottom labels
  \node[blue!70,  font=\tiny\bfseries] at (-4.8,-4.72)
    {Prior work (extended)};
  \node[orange!80,font=\tiny\bfseries] at ( 0.0,-4.72)
    {Prior work (extended)};
  \node[red!75,   font=\tiny\bfseries] at ( 4.8,-4.72)
    {New type (empirically grounded)};
\end{tikzpicture}
\caption{Three-type RS-MLLM hallucination taxonomy. Types~I and~II
extend the Aerial Mirage~\cite{aerialmirage2026} two-type framework with
refined causal analysis. Type~III (Spectral Hallucination) is new and
is empirically grounded in HM-Bench~\cite{hmbench2026}, which evaluates
18 MLLMs on HSI-specific tasks and finds systematic failure on
spatial-spectral reasoning across all tested models.}\label{fig:hall_tree}
\end{figure}

\subsubsection{Type~I: Localization Hallucination}\label{sec7:type1}

\textbf{Definition.}
The model generates a confident description of an object at a location
where that object does not exist, caused by failure to correctly localise
the query target in the RS scene.

\textbf{Causal mechanism.}
RS images cover large geographic extents: a single VHR tile at 0.3\,m
GSD spanning $2\!\times\!2$\,km contains tens of thousands of objects,
roads, and structures, presenting the model's cross-attention with an
overwhelming localisation challenge. CLIP-trained encoders exacerbate
this: the nadir-view representation of a ship from 500\,m altitude shares
almost no visual features with the ground-level ship images that dominate
CLIP training data, eliminating the visual priors that would normally
guide attention toward target-consistent regions. The result is diffuse
attention---documented by Aerial Mirage~\cite{aerialmirage2026}---in which
the model generates objects that are contextually plausible (harbour
$\to$ ships) but physically absent.

\textbf{Affected tasks.}
VQA (object existence queries), visual grounding, temporal change
detection (localising which regions changed between acquisitions).

\textbf{Distinctive example.}
A model queried about aircraft at a military airbase attends to concrete
apron markings and fuel lines---contextually associated with aircraft---and
reports the count and types of aircraft that would typically occupy such
positions, even when the apron is empty. The response is
context-conditionally hallucinated: it reflects learned co-occurrence
statistics, not the visual content of the specific image.

\textbf{Relationship to prior work.}
Directly corresponds to Aerial Mirage Type~1~\cite{aerialmirage2026} and
Seeing Clearly Type~1~\cite{seeingclearly2026}. We refine the causal
analysis by linking diffuse attention to the nadir-view domain gap
(Sect.~\ref{sec2:nadirview}) rather than treating it as a general attention
failure.

\subsubsection{Type~II: Discrimination Hallucination}\label{sec7:type2}

\textbf{Definition.}
The model correctly localises the queried region but assigns an incorrect
fine-grained class, producing a spatially accurate but semantically wrong
response.

\textbf{Causal mechanism.}
RS objects at sub-metre resolution---cars, trucks, buses, aircraft types,
ship classes---are visually similar from nadir. A compact car and a
light van at 0.5\,m GSD differ by $\approx\!3$--5 visual tokens at
ViT-L/14 patch size; the fine-grained visual information needed to resolve
this distinction is discarded by MLP or Q-Former connectors that compress
the visual token sequence. Additionally, RS instruction datasets are
heavily biased toward coarse-category labels: GeoChat-Instruct~\cite{geochat2024}
uses categories from DOTA~\cite{xia2018dota} (large-vehicle, small-vehicle,
ship) rather than sub-type labels, biasing trained models toward coarse
discrimination rather than the fine-grained classification required for
operational use.

\textbf{Evidence.}
RSHallu~\cite{rshallu2026} documents high confusion rates among similar
RS categories: car\,/\,truck\,/\,bus in parking lots; roof-type
misidentification in urban scenes; agricultural crop-type confusion.
DDFAV~\cite{ddfav2025} confirms that models trained on template-generated
instruction data systematically fail at fine-grained tasks while achieving
adequate coarse VQA performance.

\textbf{Affected tasks.}
Fine-grained VQA, OBB sub-type classification (e.g., civilian versus
military aircraft), land-cover classification at detailed category levels,
material identification.

\textbf{Relationship to prior work.}
Corresponds to Aerial Mirage Type~2~\cite{aerialmirage2026} and RSHallu
discrimination errors~\cite{rshallu2026}. We contribute additional causal
analysis connecting the failure to instruction data category distribution
and token budget allocation.

\subsubsection{Type~III: Spectral Hallucination}\label{sec7:type3}

\textbf{Definition.}
The model applies RGB visual priors to non-RGB sensor data (SAR
backscatter, HSI false-colour composites, multispectral band
combinations), generating descriptions driven by colour appearance
rather than sensor-specific physics.

\textbf{Causal mechanism (hypothesised).}
Most currently documented RS-MLLM encoders---including
CLIP~\cite{clip2021}, EVA-CLIP~\cite{fang2023evaclip},
DINOv2~\cite{oquab2023dinov2}, and RemoteCLIP~\cite{remoteclip2024}---were
pre-trained predominantly on optical RGB imagery.
While recent systems such as Earth-OneVision~\cite{earthonevision2026} extend
encoder inputs to SAR, infrared, and temporal modalities, no published RS-MLLM
encoder employs physics-specific SAR backscatter or HSI spectral pre-training
objectives. Consequently, when these encoders receive a SAR backscatter image or
an HSI false-colour composite, they may lack the sensor-physics grounding needed
to distinguish the input from a standard optical photograph. The proposed
hypothesis is that such encoders tend to map pixel values toward the nearest
RGB-trained feature cluster, without representing the physical meaning of the
sensor modality.

A concrete illustrative scenario for HSI: researchers routinely create HSI
false-colour composites by mapping the near-infrared (NIR) band to the
red display channel. In such composites, healthy vegetation appears
\emph{bright red} (high NIR reflectance). An RS-MLLM encoder, receiving
a bright-red region, may activate RGB-trained fire, bare-earth, or
red-soil features. The model could consequently describe healthy farmland as
``fire-damaged land'' or ``exposed red soil''---a response that is
confident, specific, and physically wrong. For SAR: calm water returns
nearly zero backscatter (dark pixels); double-bounce scattering from
building corners produces very bright pixels. An RGB-oriented encoder
may interpret this dark-adjacent-to-bright pattern as shadows adjacent to
illuminated structures---activating road, building, and shadow features
rather than the correct SAR scattering interpretation.

These failure patterns, if consistently reproduced, would constitute
\emph{systematic} failures: the same physically wrong interpretation reproduced
whenever the same sensor type is presented. Note that this causal account
is a mechanistic hypothesis, not a proven mechanism; controlled ablation
studies with physics-specific encoders would be needed to confirm it.

\textbf{Empirical grounding.}
HM-Bench~\cite{hmbench2026} provides controlled empirical
evidence that is consistent with the Type~III hypothesis.
The benchmark evaluates 18 representative MLLMs on 13 HSI-specific task
categories including spectral unmixing, material identification, and
multi-temporal change detection, using PCA-based false-colour composite
images and structured textual reports (rather than raw HSI cubes).
All 18 evaluated models exhibit systematic difficulty on complex
spatial--spectral reasoning tasks~\cite{hmbench2026}, providing
empirical motivation for the proposed Type~III spectral-failure
category. This pattern is consistent across all model families and
all task categories evaluated in HM-Bench.
The HM-Bench results are \emph{consistent with the hypothesis} that
RGB-oriented representations and inadequate spectral grounding contribute
to Type~III failures; however, because HM-Bench uses PCA-compressed
composites rather than raw spectral cubes, it does not directly
isolate the encoder mechanism. Controlled ablation studies comparing
RGB-trained encoders with physics-specific spectral encoders would be
needed to definitively confirm the causal account.

Supporting evidence: work on MLLM robustness to image
degradation~\cite{degradation2025} demonstrates that MLLMs trained
on RGB exhibit strong blur bias, mode collapse, and taxonomy
misalignment when presented with inputs that deviate from the RGB
distribution. Non-RGB sensor imagery represents an extreme case of
this distributional shift.

\textbf{Affected tasks.}
HSI land-cover classification, SAR scene interpretation, material
identification, crop-type mapping from multispectral imagery with
non-standard band combinations, SAR-based vessel detection.

\textbf{Why this type has not been previously identified.}
Aerial Mirage~\cite{aerialmirage2026}, RSHallu~\cite{rshallu2026},
and Seeing Clearly~\cite{seeingclearly2026} evaluate RS-MLLMs
exclusively on optical RGB imagery. DDFAV~\cite{ddfav2025} similarly
constructs its instruction data from optical sources. To the best of
our knowledge, none of these works explicitly categorises sensor-modality
misinterpretation as a distinct hallucination type.
HM-Bench~\cite{hmbench2026} is the first to systematically test non-RGB
RS inputs, revealing a failure mode that is categorically different from Types~I
and~II in both cause and mitigation strategy.

Table~\ref{tab:hall_taxonomy} summarises the three types across all
defining dimensions.

\begin{table*}[t]
\caption{Three-type RS-MLLM hallucination taxonomy. Types~I and~II extend the
Aerial Mirage~\cite{aerialmirage2026} two-type framework; Type~III is
\textbf{introduced in this survey} and is empirically grounded in
HM-Bench~\cite{hmbench2026}. \textbf{Freq.}: approximate share of RS-MLLM
hallucination failures attributable to each type, estimated from
RSHallu~\cite{rshallu2026} and HM-Bench~\cite{hmbench2026} error analyses
(optical-only models; SAR/HSI errors excluded from denominator for Types~I--II).
\textbf{Mitig.\,Diff.} (1\,=\,easy\,$\to$\,5\,=\,hardest): expert assessment
based on available mitigation literature.}\label{tab:hall_taxonomy}
\tiny\setlength\tabcolsep{2pt}
\begin{tabular*}{\textwidth}{@{\extracolsep\fill}p{1.2cm}p{2.1cm}p{2.0cm}p{1.8cm}p{1.5cm}p{2.0cm}p{2.3cm}}
\toprule
\textbf{Type} &
\textbf{Trigger} &
\textbf{Root Cause} &
\textbf{Tasks} &
\textbf{Detection} &
\textbf{Mitig.\,Diff.} &
\textbf{Real-World Risk} \\
\midrule
Type~I\\Localization
  & Target absent; attention disperses to plausible context
  & Nadir-view gap; CLIP blind-pair; scene clutter
  & VQA (exist.), grounding, change det.
  & RSHallu~\cite{rshallu2026}; Aerial Mirage~\cite{aerialmirage2026}
  & 2/5 — data-side + inference-side fixes exist~\cite{ddfav2025,seeingclearly2026}
  & Rescue teams misdirected; empty berths reported occupied \\[4pt]
Type~II\\Discrimination
  & Target localised; wrong fine-grained class
  & Token budget too small; coarse instruction data
  & Fine-grained VQA, OBB sub-type, land-cover
  & RSHallu~\cite{rshallu2026}; GeoHallu~\cite{geohallu2026}
  & 3/5 — needs higher-quality data and larger token budget
  & Structural damage misclassified; wrong crop-type payout \\[4pt]
Type~III\\Spectral \textbf{[NEW]}
  & Non-RGB input (SAR/HSI/MSI); model uses RGB priors
  & All encoders RGB-only; zero spectral-physics encoding
  & HSI classif., SAR interp., material ID
  & HM-Bench~\cite{hmbench2026} (partial; see text)
  & 5/5 — needs modality-specific encoder redesign; no fix published
  & Drought crops called healthy; SAR ship = ocean clutter \\
\botrule
\end{tabular*}
\footnotetext{Mitig.\,Diff.\,=\,Mitigation Difficulty (1\,=\,easiest; 5\,=\,hardest),
expert assessment from available mitigation literature.
[NEW]\,=\,introduced in this survey.
Frequency estimates (55\,/\,30\,/\,15\,\%) are omitted: they derive from sources
with heterogeneous denominators and cannot be combined into a single consistent measure;
the taxonomy is defensible without them.}
\end{table*}

%% ----------------------------------------------------------
\subsection{Mitigation Strategies}\label{sec7:mitig}
%% ----------------------------------------------------------

The three hallucination types require distinct mitigation strategies,
reflecting their different causal mechanisms.

\textbf{Data-side mitigation (Types~I and~II).}
DDFAV~\cite{ddfav2025} demonstrates that replacing shallow template-generated
instruction pairs with reasoning-intensive pairs substantially reduces
both Type~I and Type~II hallucination rates. The critical design principles
are: (i)~explicit negative examples---instructions querying absent objects,
ensuring the model learns to say ``no object present'' rather than generating
a plausible but wrong description; (ii)~graduated difficulty from coarse
scene-level to fine-grained object-level queries; and (iii)~diversity of
phrasing for the same spatial relationship. The DDFAV instruction corpus
(1.6\,GB) represents the largest publicly available quality-controlled RS
instruction dataset.

\textbf{Training-side mitigation (Types~I and~II).}
RemoteShield~\cite{remoteshield2026} applies cross-condition preference
learning---analogous to RLHF---to RS-MLLMs. Paired (correct, hallucinated)
response examples for RS scenes are used to train the model to prefer
physically grounded responses via a contrastive objective. This approach
directly targets the model's tendency to generate contextually plausible
but visually ungrounded responses.

\textbf{Inference-side mitigation (Type~I).}
Seeing Clearly~\cite{seeingclearly2026} proposes a training-free
inference-time correction that identifies diffuse attention patterns
(the signature of Type~I) and applies attention sharpening before
response generation. The method requires no additional fine-tuning and
is applicable to deployed RS-MLLMs. Its limitation is that it operates
on attention maps and cannot address the pre-attentional Type~III
failure: spectral hallucination occurs in the encoder before
cross-attention is computed.

\textbf{Type~III mitigation (future direction).}
No published method directly addresses spectral hallucination. Three
research directions are plausible. First, modality-conditioned encoders:
pre-train separate encoder branches for SAR and HSI with
sensor-physics-informed objectives (predicting SAR backscatter
coefficient distributions; predicting spectral unmixing endmembers)
rather than contrastive image-text objectives. Second,
physics-informed loss terms: augment MLLM training with a spectral
consistency loss that penalises responses inconsistent with known
sensor-physics constraints. The Rayleigh-constrained SAR-MSI fusion
approach in our prior work~\cite{rayleigh2026} demonstrates that
physics constraints improve cross-sensor feature alignment; an
analogous constraint applied to the MLLM training objective could
suppress Type~III hallucination. Third, modality prompt tokens:
explicitly signal to the LLM backbone which sensor modality is being
processed, analogous to language tokens in multilingual LLMs, enabling
the model to switch its interpretation context.

%% ----------------------------------------------------------
\subsection{Hallucination Benchmarks and Gaps}\label{sec7:bench}
%% ----------------------------------------------------------

Table~\ref{tab:hall_bench} compares the four available RS hallucination
benchmarks across the dimensions of our taxonomy. Values are as reported
in the respective papers.

\begin{table*}[t]
\caption{RS-MLLM hallucination benchmark coverage. ``\checkmark''\,=\,explicitly
evaluated; ``(\checkmark)''\,=\,partially evaluated; ``$\times$''\,=\,not evaluated.
All benchmark characteristics as reported in the respective
papers.}\label{tab:hall_bench}
\footnotesize\setlength\tabcolsep{3pt}
\begin{tabular*}{\textwidth}{@{\extracolsep\fill}lllcccllc}
\toprule
\textbf{Benchmark} & \textbf{Year} & \textbf{Venue} &
\textbf{Type~I} & \textbf{Type~II} & \textbf{Type~III} &
\textbf{Scale} & \textbf{Eval.} & \textbf{Modal.} \\
\midrule
DDFAV~\cite{ddfav2025}
  & 2025 & MDPI RS
  & $(\checkmark)$ & $(\checkmark)$ & $\times$
  & 1.6\,GB & VQA accuracy & Optical \\[3pt]
RSHallu~\cite{rshallu2026}
  & 2026 & [Pre]
  & \checkmark & \checkmark & $\times$
  & NR & Exist.\,/\,attr.\,/\,rel. & Optical \\[3pt]
Aerial Mirage~\cite{aerialmirage2026}
  & 2026 & [Pre]
  & \checkmark & \checkmark & $\times$
  & NR & 2-type classif. & Optical \\[3pt]
HM-Bench~\cite{hmbench2026}
  & 2026 & [Pre]
  & $\times$ & $(\checkmark)$ & $(\checkmark)$
  & 19,337 & 13 task cats. & HSI \\
\botrule
\end{tabular*}
\footnotetext{NR\,=\,Not reported in the paper.
HM-Bench~\cite{hmbench2026} evaluates task performance on HSI tasks
rather than per-query hallucination labels; its Type~III coverage is
therefore partial. No benchmark provides comprehensive evaluation of
Type~III (Spectral) hallucination. DDFAV, RSHallu, and Aerial Mirage
evaluate exclusively on optical RGB imagery.}
\end{table*}

Three benchmark gaps are evident from Table~\ref{tab:hall_bench}. First,
all optical benchmarks (DDFAV, RSHallu, Aerial Mirage) evaluate
exclusively on optical RGB imagery, leaving Type~III entirely untested
and providing no SAR-specific evaluation. Second,
HM-Bench~\cite{hmbench2026} tests non-optical modalities (HSI) but
reports task performance rather than per-query hallucination
labels---making it a \emph{partial} Type~III resource rather than a
dedicated benchmark. Third, no existing benchmark evaluates cross-type
hallucination, in which a single query triggers failures at multiple
stages simultaneously (e.g., a Type~I localisation failure followed by
a Type~III spectral misinterpretation of the wrongly attended region).

\paragraph{The missing Type~III benchmark.}
A comprehensive Type~III evaluation benchmark requires: (i)~co-registered
optical, SAR, and HSI images of the same scene; (ii)~per-modality
query--response pairs with ground-truth labels independent of colour
appearance; and (iii)~annotation of whether the model's response was
driven by spectral physics or RGB visual priors. No such benchmark
exists. Its construction would require approximately 5,000--10,000
co-registered triplets with expert annotation---a tractable scale that
represents a concrete and high-impact research contribution for the
RS-MLLM community.

\subsection*{Summary}

This section introduced a systematic three-type hallucination
taxonomy for RS-MLLMs that proposes a formal three-way partition
(Localization, Discrimination, and Spectral failure modes) with distinct
causal mechanisms grounded in the RS domain challenges of Sect.~\ref{sec2}.
To the best of our knowledge, this specific partition with its RS-physics
causal grounding has not been proposed in prior RS-MLLM hallucination work.
Type~I (Localization) and Type~II (Discrimination)
extend and refine the Aerial Mirage~\cite{aerialmirage2026} framework with
additional causal analysis rooted in the RS domain challenges of
Sect.~\ref{sec2}. Type~III (Spectral Hallucination) is a new type,
empirically grounded in HM-Bench~\cite{hmbench2026}, that captures a
failure mode not, to the best of our knowledge, previously isolated as a
distinct category in the RS hallucination literature; our taxonomy extends
prior localisation/discrimination frameworks by incorporating this spectral
failure evidence. This type requires modality-specific encoder design for
mitigation. The benchmark analysis (Table~\ref{tab:hall_bench}) confirms that
Type~III remains unaddressed by any existing evaluation framework, establishing
a clear research priority for the RS-MLLM community.
